\documentclass[10pt, a4paper, landscape]{article}

\usepackage[T1]{fontenc}
\usepackage[utf8]{inputenc}
\usepackage{tgpagella}
\usepackage{tgcursor}
\usepackage{microtype}
\usepackage[spanish, es-noshorthands]{babel}
\usepackage[table, xcdraw]{xcolor}
\usepackage{booktabs}
\usepackage{tabularx}
\usepackage{array}
\usepackage{amsmath, amssymb}
\usepackage[most]{tcolorbox}
\usepackage{tikz}
\usetikzlibrary{shapes.geometric, arrows.meta, positioning, calc}
\usepackage[margin=1.4cm]{geometry}
\usepackage{fancyhdr}

\pagestyle{fancy}
\fancyhf{}
\setlength{\headheight}{16pt}
\lhead{\textbf{IMT-342}}
\chead{Hoja Docente: Auditoría y Cierre Operativo de Hito 2}
\rhead{Semana 12 - Sesión 2}
\lfoot{Control de Defensa Individual y Condición para Recuperatorio W13}
\rfoot{Página \thepage}
\renewcommand{\headrulewidth}{0.4pt}
\renewcommand{\footrulewidth}{0.4pt}

\definecolor{ucbblue}{RGB}{0, 51, 102}

\begin{document}

\begin{center}
    \Large\textbf{\textcolor{ucbblue}{Hoja de Control Docente: Cierre Operativo de Hito 2 y Dependencia del Recuperatorio W13}}[cite: 16]
\end{center}

\section*{1. Flujo de Decisión Operativa en el Segmento 5 (Visual 6)[cite: 16]}
El Hito 2 \textbf{no} se fusiona con el Hito 3 ni traslada su nota[cite: 16]. Su cierre formal (con entregables validados y evidencia de defensa individual conforme a la rúbrica original) junto a la calificación de EC1/EC2 es requisito previo para consolidar las Unidades I y II y convocar el recuperatorio de W13[cite: 16].

\begin{center}
% Visual 6: Gestión operativa Hito 2
\begin{tikzpicture}[
    box/.style={rectangle, draw=ucbblue, thick, fill=ucbblue!5, rounded corners, text width=4.6cm, align=center, font=\scriptsize, minimum height=1.1cm},
    decbox/.style={rectangle, draw=orange!80!black, thick, fill=yellow!10, rounded corners, text width=5.2cm, align=center, font=\scriptsize, minimum height=1.1cm},
    recbox/.style={rectangle, draw=green!60!black, thick, fill=green!10, rounded corners, text width=4.8cm, align=center, font=\scriptsize, minimum height=1.1cm},
    arrow/.style={-Stealth, very thick, ucbblue}
]
    \node[box] (e1) at (0, 0) {\textbf{1. Verificación de Entregables}\\$\square$ \texttt{.ipynb} ejecutable verificado\\$\square$ Informe PDF IEEE recibido[cite: 16]};
    \node[box] (e2) at (5.8, 0) {\textbf{2. Auditoría de Módulos Unidad II}\\Pieper, FK, LSPB, Jacobiano, Singularidad, Dinámica 1R[cite: 16]};
    \node[decbox] (e3) at (11.9, 0) {\textbf{3. Decisión al Cierre del Segmento 5}\\(a) Defensa hoy con evidencia individual, o\\(b) Convocatoria externa fechada, o\\(c) Bloqueo concreto + responsable[cite: 16]};
    \node[recbox] (e4) at (18.2, 0) {\textbf{4. Consolidación Unidades I–II}\\Cierre Hito 2 + Corrección EC1/EC2\\$\implies$ Habilita Recuperatorio W13[cite: 16]};

    \draw[arrow] (e1) -- (e2);
    \draw[arrow] (e2) -- (e3);
    \draw[arrow] (e3) -- (e4);
\end{tikzpicture}
\end{center}

\section*{2. Matriz de Registro por Grupo (Segmento 5 — 20 min)[cite: 16]}
\begin{center}
\scriptsize
\setlength{\tabcolsep}{3pt}
\renewcommand{\arraystretch}{1.55}
% Suma de \hsize = 0.45 + 1.15 + 0.70 + 0.95 + 0.70 + 0.75 + 0.95 + 0.65 + 1.10 + 1.60 = 9.0 para las 9 columnas X (+ 1 col c)
\begin{tabularx}{\textwidth}{|c|>{\hsize=1.15\hsize\raggedright\arraybackslash}X|>{\hsize=0.75\hsize\centering\arraybackslash}X|>{\hsize=0.95\hsize\centering\arraybackslash}X|>{\hsize=0.75\hsize\centering\arraybackslash}X|>{\hsize=0.75\hsize\centering\arraybackslash}X|>{\hsize=1.00\hsize\centering\arraybackslash}X|>{\hsize=0.70\hsize\centering\arraybackslash}X|>{\hsize=1.20\hsize\raggedright\arraybackslash}X|>{\hsize=1.75\hsize\raggedright\arraybackslash}X|}
\hline
\rowcolor{ucbblue!15}
\textbf{Grp.}[cite: 16] & \textbf{Integrantes}[cite: 16] & \textbf{\texttt{.ipynb} Recibido}[cite: 16] & \textbf{Reproducibilidad} (\texttt{verif. / no verif. / falló})[cite: 16] & \textbf{IEEE PDF Recibido}[cite: 16] & \textbf{Listo para Defensa}[cite: 16] & \textbf{Evidencia Indiv. Registrada}[cite: 16] & \textbf{Duración Acordada}[cite: 16] & \textbf{Día / Hora de Convocatoria}[cite: 16] & \textbf{Bloqueo Identificado + Responsable y Fecha Límite}[cite: 16] \\ \hline
G1 & & sí / no & & sí / no & sí / no & sí / no & & & \\ \hline
G2 & & sí / no & & sí / no & sí / no & sí / no & & & \\ \hline
G3 & & sí / no & & sí / no & sí / no & sí / no & & & \\ \hline
G4 & & sí / no & & sí / no & sí / no & sí / no & & & \\ \hline
G5 & & sí / no & & sí / no & sí / no & sí / no & & & \\ \hline
G6 & & sí / no & & sí / no & sí / no & sí / no & & & \\ \hline
\end{tabularx}
\end{center}

\vspace{0.2cm}
\begin{tcolorbox}[colback=red!5, colframe=red!70!black, title=\textbf{Regla de Integridad Evaluativa para Hito 2} {[cite: 16]}, left=2mm, right=2mm, top=1mm, bottom=1mm]
    \scriptsize\raggedright
    No confundir \textit{«se entregó el notebook»} con \textit{«cada integrante demostró dominio individual»}[cite: 16]. Tampoco asumir que el archivo PDF acredita una fuente LaTeX si no está documentada[cite: 16]. La microdefensa dentro de la sesión solo procede si se registra evidencia individual rubricada sin comprometer la supervisión del aula/laboratorio; en caso contrario, debe quedar firmada la fecha de convocatoria externa o el bloqueo concreto[cite: 16].
\end{tcolorbox}

\end{document}
